\section{敏感性、边界与验证}
\subsection{窗口参数的可辨识性}
候选规则全部一致应用于六个图，不按算子名称选用专门参数。为分离窗口影响，表\ref{tab:sensitivity}只比较释放权重 $\alpha=1$、距离换出、最佳适配都固定时的 $W=32$ 和 $W=128$。$H_W$ 是原始序列峰值，$D_W$ 是随后分配的额外搬运量。
\begin{table}[H]\centering\caption{固定其他规则时的窗口对照}\label{tab:sensitivity}\zihao{-4}
\begin{tabular}{lrrrrr}\toprule
计算图 & $H_{32}$ & $H_{128}$ & $D_{32}$ & $D_{128}$ & 顺序相同\\\midrule
卷积小图 & 21762 & 21762 & 95564 & 95564 & 否\\
卷积大图 & 157630 & 180030 & 407112 & 439760 & 否\\
注意力小图 & 5536 & 5536 & 24896 & 24896 & 是\\
注意力大图 & 19008 & 19008 & 129512 & 129512 & 否\\
矩阵乘小图 & 14208 & 24704 & 85632 & 152832 & 否\\
矩阵乘大图 & 39680 & 52352 & 611968 & 719488 & 否\\
\bottomrule\end{tabular}
\end{table}
矩阵乘小图扩窗后，峰值从 14208 升至 24704，搬运量从 85632 升至 152832，说明扩大就绪选择范围不会单调改善结果。窗口改变的是每一步的可选集合，不是对已完成全局方案作包含关系扩展；一旦局部选择改变，后续状态也改变，所以“更多选择”不保证最终更优。

卷积小图与注意力大图两种窗口的峰值和搬运量相同，但完整序列并不相同，时间分别为 642017 与 640953、485807 与 486988 周期。这是目标统计量无法区分某些顺序差异的例子，不能据此认定参数未生效。注意力小图的两份完整序列确实相同，只能认定这两个窗口值在该实例上落在相同输出区间；本轮未记录每次候选竞争，尚不能区分窗口未约束就绪集合与最优局部选择恰好一致这两种原因。

此外，$(64,2)$ 同时改变窗口和权重，因此只作为另一个算法候选，不把其与 $(32,1)$ 的差异单独归因于释放权重。本轮未进行容量扰动实验：改变容量意味着改变题目边界条件，现有结果不足以支持对任意缓存规模的稳健性结论。

\subsection{手算对照与错误注入}
可行性检查从导出的序列、偏移与换出列表重新读入，不调用候选生成或分配函数。它检查原始及新增节点恰好出现一次、原始边顺序、换出配对、每次地址申请的整数边界与逐单元占用、操作使用时的驻留状态，最后重新计算流水和搬运量。问题三与问题二的十二份完整方案及六份问题一序列均完成相应重算。

为了防止执行时间计算把所有操作错误串行化，构造两个无依赖操作，分别在输入、输出搬运单元执行 10 和 20 周期，缓冲区大小均为 4。地址互不相交时，正确总时间为 $\max(10,20)=20$；复用同一地址时，前段释放到后段申请形成依赖，正确总时间为 $10+20=30$。这两个答案只差地址关系，可同时检验并行能力和地址等待是否生效。

对大小为 4 的输入复制缓冲区，一次换出搬运量为 4，换入耗时为 $2\times4+150=158$；按测试给定的其他操作和依赖，完整总时间为 188。对应非输入复制对象，搬运量为 8，两次搬运各耗时 158，完整时间为 346。这些手算测试覆盖零耗时换出和双向计费。再分别注入地址重叠、拓扑逆序、越界偏移及节点重复，检查程序均拒绝对应附件。拒绝非法输入只是验证的一部分，不替代对合法实例的数学建模判断。

\subsection{解释边界与尚未解决的问题}
本题附录 C 的地址复用描述直接针对 ALLOC/FREE；本文将其延伸为驻留段结束到开始的依赖。若把已换出的缓冲区仍视为占用地址直到最终 FREE，就无法表达换出释放空间的目的，并可能生成环。因此采用驻留段解释有物理上的一致性，但尚无官方评价器作交叉确认。相关时间结论均以此解释为前提。

对搬运量，本文只有非负这一平凡下界，没有得到六图的最小必要换出集合；对驻留峰值，式\eqref{eq:hlower}也较弱。这使模型能证明方案满足约束，却不能证明搬运和峰值接近全局最优。时间下界更有辨识力，但强下界仅针对固定事件或驻留图。三种解质量保证的范围不同，不应合并成一个“算法最优”结论。

六图都是题面示例，没有独立保留的泛化测试集；相同规则跨图应用，只证明实现没有写入算子专用分支。进一步的通用性评价需要不同规模和不同共享数据结构的新图。求解与重放虽分开实现，仍共享输入解析和题意解释，不能排除共同理解偏差。真实芯片测量也未进行，周期改善不等同于实测吞吐加速。
